% Review
\subsection{Scrimba}

Posledním analyzovaným řešením je webová aplikace Scrimba. Tato aplikace se od ostatních odlišuje zejména svým inovativním způsobem výuky programování, který kombinuje prostředí připomínající skutečné IDE a interaktivní výuková videa. Lekce se tak odehrávají přímo v editoru kódu, ve kterém je uživateli vysvětlována látka a zadávány úkoly. Tento přístup výuky může přinést řadu výhod oproti tradičním metodám výuky a proto jsem se rozhodl tuto aplikaci do analýzy zahrnout.

Platforma Scrimba obsahuje přes 60 kurzů na různá témata napříč softwarovým inženýrstvím. V době psaní tohoto textu jsou kurzy zaměřeny zejména na webové technologie a umělou inteligenci. 

Na rozdíl od ostatních analyzovaných aplikací Scrimba nemá v současné době mobilní verzi aplikace. Proto pro analýzu využijí pouze webovou aplikaci Scrimba. 

% Review
\subsubsection{Vzhled a uživatelská zkušenost}

% Odstavec: Obecný dojem (komplexita/přehlednost UI, paleta barev, typografie)

Na první pohled působí Scrimba jako přehledná a minimalistická aplikace. Ve srovnání s konkurenční aplikací Mimo používá aplikace menší velikost písma a zobrazuje na stránkách více textu. Aplikace tak může působit komplexnějším dojmem. Co se palety barev týče, používá Scrimba monochromatickou paletu modré barvy. Odstíny modré jsou použity pro pozadí, primární akce, ikony a textové prvky. Pro upozornění a některé štítky aplikace vyjímečně využívá i žlutou nebo zelenou barvu. Napříč aplikací je používáno systémové písmo, tedy písmo, které je nastaveno v závislosti na operačním systému uživatele. Na systému Windows je tak použito například písmo Segoe, na zařízeních macOS písmo San Francisco a na zařízeních Android písmo Roboto.  % Write - přidat referenci sem https://fonts.google.com/knowledge/glossary/system_font_web_safe_font

% Odstavec: Rozložení prvků na obrazovce (lekce v kurzu, učení)

Prvky jsou na obrazovce rozloženy velice přehledně a logicky. Ve srovnání s ostatními aplikacemi Scrimba používá méně prostoru mezi jednotlivými prvky uživatelského rozhraní. V důsledku toho může aplikace místy působit komplexnějším dojmem.

Na hlavní stránce kurzu jsou jednotlivé sekce, moduly a lekce vizuálně uspořádány do stromové struktury, kterou lze postupně rozbalovat. Díky tomu může uživatel jednoduše získat přehled o celém kurzu a rozbalit části, které ho zajímají. Lekce jsou tak uspořádány velice přehledně.

% Odstavec: Navigace v aplikaci

Na rozdíl od ostatních analyzovaných aplikací Scrimba nevyužívá pro navigaci horní navigační lištu, ani žádnou patičku stránky. Pro navigaci je využívána výhradně boční navigační lišta, která je umístěna na levé straně obrazovky.

% Odstavec: Animace a efekty a mikro interakce
 
Ve srovnání s ostatními aplikacemi Scrimba výrazně více využívá animace jednotlivých prvků uživatelského rozhraní a animované přechody mezi stránkami. Na většině stránek je při jejich načtení některý z prvků animován. Další výraznou animaci lze zpozorovat při přejetí kurzorem nad tlačítko. Při animaci je tlačítko zvýrazněno čarou, která ho postupně obklopí po jeho obvodu.

Animace jsou napříč aplikací používány často, nicméně minimalisticky a v přiměřeném množství. Díky tomu Scrimba může působit výrazně interaktivnějším a vizuálně poutavějším dojmem, než její konkurenti.

% Review
\subsubsection{Způsoby učení}

% Struktura kurzu (lineární, tematické moduly, career paths, kapitoly atd.)
Scrimba nabízí 

% Jak se spustí učení a jak vypadá study session
Uživateli je při učení v prohlížeči zobrazen editor připomínající skutečné IDE, ve kterém uživatel může prohlížet soubory a editovat kód. Lekce jsou strukturovány tak, že uživatel poslouchá audio s vysvětlením dané látky a sleduje přímo v editoru, jak lektor přepisuje kód a jak pohybuje kurzorem.

% Druhy cvičení

% Review
\subsubsection{Způsoby motivace uživatele}

% Jaké formy gamifikace aplikace používá


% Review
\subsubsection{Způsoby monetizace}

% Free tier

% Premium tier


% Free tier


% Review
\subsubsection{Silné a slabé stránky}

% Silné stránky


% Slabé stránky


% Review
\subsubsection{Další poznatky}

% Zajímavé funkce

% Write - kurzy scrimby a některé lekce lze prohlížet i bez přihlášení. Díky tomu může uživatel 

% (Zajímavé detaily nebo cokoliv co mě napadne)
